\chapter{模板兼容性检查}
\label{template-chapter-check}

\section{中文、公式和 Stacks 数学宏}
\label{template-section-macros}

中文正文、行内公式 $X \longrightarrow Y$ 和常用算子应能混合排版。例如
$\Spec(R)$、$\Hom_R(M,N)$、$\SheafExt^1(\mathcal F,\mathcal G)$、
$\QCohstack$、$\etale$ 与 $\colim_i M_i$ 均来自英文 harvest 的共享前导宏。

\begin{definition}[排版测试]
\label{template-definition-typesetting}
这是一个用于推进共享 \texttt{subsection} 计数器的定义环境。
\end{definition}

\begin{equation}
\label{template-equation-polynomial}
  f(x)=x^2+1.
\end{equation}

公式~\ref{template-equation-polynomial}、定义~\ref{template-definition-typesetting}
以及本章第~\ref{template-section-diagrams}~节的交叉引用都应正确。

标准 \texttt{matrix} 环境和粗体矩阵命令不能冲突：
\[
  \begin{matrix}
    a & b \\
    c & d
  \end{matrix}
  \qquad \mat{A}.
\]

参照入口中的数学宏包和辅助命令也应可直接使用：
\[
  \diag{1,2},\qquad
  \closure{U},\qquad
  \seg{0}{1},\ \segl{0}{1},\ \segr{0}{1},\ \seglr{0}{1},\qquad
  \dv{f}{x}.
\]
量纲命令 \SI{3}{\metre}、小写罗马数字 \romannum{4}，以及
\texttt{NiceMatrix} 环境均应正常：
\[
  \begin{NiceMatrix}
    1 & 0 \\
    0 & 1
  \end{NiceMatrix}.
\]

\section{定理、证明和编辑元数据}
\label{template-section-theorems}

\begin{theorem}[示例定理]
\label{template-theorem-example}
定理标题、编号和中文正文应当正常显示。
\end{theorem}

\begin{proof}
证明环境应使用中文标题，并在结尾显示证毕符号。
\end{proof}

\begin{lemma}
\label{template-lemma-example}
引理与定理共用 Stacks Project 的 \texttt{subsection} 计数器。
\end{lemma}

\begin{proposition}
命题环境可用。
\end{proposition}

\begin{exercise}
\label{template-solution}
练习环境可用，并为可选的解答章节提供稳定标签。
\end{exercise}

\begin{situation}
\label{template-situation-example}
英文源文件特有的 \texttt{situation} 环境可用。
\end{situation}

\begin{remark}
单个注记环境可用。
\end{remark}

\begin{remarks}
复数形式的注记环境也可用。
\end{remarks}

\begin{reference}
这段编辑性参考信息不应出现在成书 PDF 中。
\end{reference}
\begin{slogan}
这段口号元数据不应出现在成书 PDF 中。
\end{slogan}
\begin{history}
这段历史元数据不应出现在成书 PDF 中。
\end{history}

\section{交换图、标签链接和索引}
\label{template-section-diagrams}

Xy-pic 交换图应能直接沿用英文源文件的写法：
\[
\xymatrix{
  X \ar[r]^f \ar[d]_g & Y \ar[d]^h \\
  Z \ar[r]_k & W
}
\]

永久标签链接示例：\StacksTag{0000}。项目引用示例见~\cite{stacks-project}。

\index{Stacks Project}
\index{中文排版}
索引应能生成，且不产生按 UTF-8 字节计算的乱码分组标题。
